\section{Introdução}

\subsection{Descrição do Algoritmo}

O K-Means é um algoritmo de agrupamento não supervisionado que particiona um
conjunto de dados em \(k\) grupos (clusters), de forma que cada ponto pertença
ao cluster cujo centróide é o mais próximo. O algoritmo opera em iterações:
na etapa de \textbf{atribuição}, cada ponto é associado ao centróide mais próximo
segundo a distância euclidiana; na etapa de \textbf{atualização}, cada centróide
é recalculado como a média aritmética de todos os pontos a ele atribuídos. Esse
ciclo se repete até que os centróides estabilizem (deslocamento máximo abaixo de
uma tolerância) ou o número máximo de iterações seja atingido.

A complexidade por iteração é \(O(N \cdot k \cdot d)\), onde \(N\) é o número de
pontos, \(k\) o número de clusters e \(d\) a dimensionalidade dos dados. Por
percorrer todos os pontos a cada iteração, o algoritmo é naturalmente paralelizável
na etapa de atribuição, tornando-o um caso de estudo relevante para comparação
entre estratégias de concorrência.

\subsection{Descrição da Linguagem Erlang}

Erlang é uma linguagem funcional criada pela Ericsson, projetada para sistemas distribuídos e
tolerantes a falhas. Executa sobre a máquina virtual BEAM, que provê concorrência nativa por
meio de processos leves e isolados.

\textbf{Modelo de programação.}
Utiliza passagem de mensagens (modelo de atores). Processos não compartilham memória e se
comunicam exclusivamente via caixas postais, eliminando por construção os problemas de acesso
concorrente a estado compartilhado.

\textbf{Modelo de memória e visibilidade.}
Cada processo possui heap privado. Mensagens são copiadas entre heaps ao ser enviadas, o que
garante isolamento total. Por não haver memória compartilhada mutável, questões de
visibilidade — como barreiras de memória e \texttt{volatile} do Java — não se aplicam.

\textbf{Controle de região crítica.}
Não existem mutexes ou semáforos. O acesso a recursos exclusivos é feito encapsulando-os em
um processo servidor que atende requisições sequencialmente — padrão formalizado pelo
\texttt{gen\_server} OTP.

\textbf{Coleta de lixo.}
O GC opera por processo, de forma independente, sem pausas globais (\textit{stop-the-world}).
Processos de vida curta têm seus heaps simplesmente descartados ao terminar.

\textbf{JIT.}
A partir do OTP 24, a BEAM inclui um compilador JIT (AsmJit) que gera código nativo para
x86-64 e AArch64 no carregamento dos módulos, com ganhos típicos de 20--40\% em relação ao
bytecode interpretado.